% Shared preamble for the Week 7 student pages (compile with xelatex).
% Each band file sets \bandname and the body font size before % Shared preamble for the Week 7 student pages (compile with xelatex).
% Each band file sets \bandname and the body font size before % Shared preamble for the Week 7 student pages (compile with xelatex).
% Each band file sets \bandname and the body font size before % Shared preamble for the Week 7 student pages (compile with xelatex).
% Each band file sets \bandname and the body font size before \input{common}.
\usepackage[letterpaper,top=0.55in,bottom=0.55in,left=0.65in,right=0.65in,
  includeheadfoot,headheight=16pt,headsep=12pt,footskip=26pt]{geometry}
\usepackage{fontspec}
\setmainfont{Andika}
\usepackage{tikz}
\usetikzlibrary{arrows.meta,calc,positioning}
\usepackage{fancyhdr}
\usepackage{array}
\usepackage{enumitem}

\pagestyle{fancy}
\fancyhf{}
\fancyhead[L]{\normalsize Week 7 / Take-away games / \bandname}
\fancyfoot[L]{\normalsize Bellingham Math Circle / Week 7}
\fancyfoot[R]{\normalsize \thepage}
\renewcommand{\headrulewidth}{0.4pt}
\renewcommand{\footrulewidth}{0pt}

\setlength{\parindent}{0pt}
\setlength{\parskip}{0pt}
\raggedright

\definecolor{counterfill}{gray}{0.80}
\definecolor{faded}{gray}{0.62}
\definecolor{lightline}{gray}{0.55}

% \prob{N} starts a problem.
\newcommand{\prob}[1]{\textbf{Problem #1:}\ }

% Ruled writing lines: \writelines{number}{spacing}
\newcommand{\writelines}[2]{%
  \par\begin{tikzpicture}
    \foreach \i in {1,...,#1} {
      \draw[lightline,line width=0.5pt] (0,-\i*#2) -- (\linewidth-1pt,-\i*#2);
    }
    \path (0,0) (0,-#1*#2-0.05);
  \end{tikzpicture}\par}

% One counter at a point: \cnt{(x,y)}{radius}
\newcommand{\cnt}[2]{\draw[line width=0.9pt,fill=counterfill] #1 circle (#2);}
% A counter that the partner already took: dashed, unfilled.
\newcommand{\gone}[2]{\draw[faded,line width=0.9pt,dash pattern=on 2.2pt off 1.8pt] #1 circle (#2);}

% Pile of n counters in rows of 5, top-left counter centre at (x0,y0),
% radius r, pitch p (inside a tikzpicture).  \pileat{n}{x0}{y0}{r}{p}
\newcommand{\pileat}[5]{%
  \foreach \k in {1,...,#1} {
    \pgfmathtruncatemacro{\pilerow}{(\k-1)/5}
    \pgfmathtruncatemacro{\pilecol}{mod(\k-1,5)}
    \pgfmathsetmacro{\pilex}{#2+\pilecol*#5}
    \pgfmathsetmacro{\piley}{#3-\pilerow*#5}
    \cnt{(\pilex,\piley)}{#4}
  }}

% Numbered boxes, numbers above: \numboxes{from}{to}{per row}{box width cm}{box height cm}
\newcommand{\numboxes}[5]{%
  \begin{tikzpicture}
    \foreach \n in {#1,...,#2} {
      \pgfmathtruncatemacro{\row}{(\n-#1)/#3}
      \pgfmathtruncatemacro{\col}{mod(\n-#1,#3)}
      \pgfmathsetmacro{\x}{\col*#4}
      \pgfmathsetmacro{\y}{-\row*(#5+0.62)}
      \node[anchor=south,inner sep=1pt] at ({\x+#4/2},{\y+0.02}) {\n};
      \draw[line width=0.7pt] (\x,\y) rectangle ({\x+#4},{\y-#5});
    }
  \end{tikzpicture}}
.
\usepackage[letterpaper,top=0.55in,bottom=0.55in,left=0.65in,right=0.65in,
  includeheadfoot,headheight=16pt,headsep=12pt,footskip=26pt]{geometry}
\usepackage{fontspec}
\setmainfont{Andika}
\usepackage{tikz}
\usetikzlibrary{arrows.meta,calc,positioning}
\usepackage{fancyhdr}
\usepackage{array}
\usepackage{enumitem}

\pagestyle{fancy}
\fancyhf{}
\fancyhead[L]{\normalsize Week 7 / Take-away games / \bandname}
\fancyfoot[L]{\normalsize Bellingham Math Circle / Week 7}
\fancyfoot[R]{\normalsize \thepage}
\renewcommand{\headrulewidth}{0.4pt}
\renewcommand{\footrulewidth}{0pt}

\setlength{\parindent}{0pt}
\setlength{\parskip}{0pt}
\raggedright

\definecolor{counterfill}{gray}{0.80}
\definecolor{faded}{gray}{0.62}
\definecolor{lightline}{gray}{0.55}

% \prob{N} starts a problem.
\newcommand{\prob}[1]{\textbf{Problem #1:}\ }

% Ruled writing lines: \writelines{number}{spacing}
\newcommand{\writelines}[2]{%
  \par\begin{tikzpicture}
    \foreach \i in {1,...,#1} {
      \draw[lightline,line width=0.5pt] (0,-\i*#2) -- (\linewidth-1pt,-\i*#2);
    }
    \path (0,0) (0,-#1*#2-0.05);
  \end{tikzpicture}\par}

% One counter at a point: \cnt{(x,y)}{radius}
\newcommand{\cnt}[2]{\draw[line width=0.9pt,fill=counterfill] #1 circle (#2);}
% A counter that the partner already took: dashed, unfilled.
\newcommand{\gone}[2]{\draw[faded,line width=0.9pt,dash pattern=on 2.2pt off 1.8pt] #1 circle (#2);}

% Pile of n counters in rows of 5, top-left counter centre at (x0,y0),
% radius r, pitch p (inside a tikzpicture).  \pileat{n}{x0}{y0}{r}{p}
\newcommand{\pileat}[5]{%
  \foreach \k in {1,...,#1} {
    \pgfmathtruncatemacro{\pilerow}{(\k-1)/5}
    \pgfmathtruncatemacro{\pilecol}{mod(\k-1,5)}
    \pgfmathsetmacro{\pilex}{#2+\pilecol*#5}
    \pgfmathsetmacro{\piley}{#3-\pilerow*#5}
    \cnt{(\pilex,\piley)}{#4}
  }}

% Numbered boxes, numbers above: \numboxes{from}{to}{per row}{box width cm}{box height cm}
\newcommand{\numboxes}[5]{%
  \begin{tikzpicture}
    \foreach \n in {#1,...,#2} {
      \pgfmathtruncatemacro{\row}{(\n-#1)/#3}
      \pgfmathtruncatemacro{\col}{mod(\n-#1,#3)}
      \pgfmathsetmacro{\x}{\col*#4}
      \pgfmathsetmacro{\y}{-\row*(#5+0.62)}
      \node[anchor=south,inner sep=1pt] at ({\x+#4/2},{\y+0.02}) {\n};
      \draw[line width=0.7pt] (\x,\y) rectangle ({\x+#4},{\y-#5});
    }
  \end{tikzpicture}}
.
\usepackage[letterpaper,top=0.55in,bottom=0.55in,left=0.65in,right=0.65in,
  includeheadfoot,headheight=16pt,headsep=12pt,footskip=26pt]{geometry}
\usepackage{fontspec}
\setmainfont{Andika}
\usepackage{tikz}
\usetikzlibrary{arrows.meta,calc,positioning}
\usepackage{fancyhdr}
\usepackage{array}
\usepackage{enumitem}

\pagestyle{fancy}
\fancyhf{}
\fancyhead[L]{\normalsize Week 7 / Take-away games / \bandname}
\fancyfoot[L]{\normalsize Bellingham Math Circle / Week 7}
\fancyfoot[R]{\normalsize \thepage}
\renewcommand{\headrulewidth}{0.4pt}
\renewcommand{\footrulewidth}{0pt}

\setlength{\parindent}{0pt}
\setlength{\parskip}{0pt}
\raggedright

\definecolor{counterfill}{gray}{0.80}
\definecolor{faded}{gray}{0.62}
\definecolor{lightline}{gray}{0.55}

% \prob{N} starts a problem.
\newcommand{\prob}[1]{\textbf{Problem #1:}\ }

% Ruled writing lines: \writelines{number}{spacing}
\newcommand{\writelines}[2]{%
  \par\begin{tikzpicture}
    \foreach \i in {1,...,#1} {
      \draw[lightline,line width=0.5pt] (0,-\i*#2) -- (\linewidth-1pt,-\i*#2);
    }
    \path (0,0) (0,-#1*#2-0.05);
  \end{tikzpicture}\par}

% One counter at a point: \cnt{(x,y)}{radius}
\newcommand{\cnt}[2]{\draw[line width=0.9pt,fill=counterfill] #1 circle (#2);}
% A counter that the partner already took: dashed, unfilled.
\newcommand{\gone}[2]{\draw[faded,line width=0.9pt,dash pattern=on 2.2pt off 1.8pt] #1 circle (#2);}

% Pile of n counters in rows of 5, top-left counter centre at (x0,y0),
% radius r, pitch p (inside a tikzpicture).  \pileat{n}{x0}{y0}{r}{p}
\newcommand{\pileat}[5]{%
  \foreach \k in {1,...,#1} {
    \pgfmathtruncatemacro{\pilerow}{(\k-1)/5}
    \pgfmathtruncatemacro{\pilecol}{mod(\k-1,5)}
    \pgfmathsetmacro{\pilex}{#2+\pilecol*#5}
    \pgfmathsetmacro{\piley}{#3-\pilerow*#5}
    \cnt{(\pilex,\piley)}{#4}
  }}

% Numbered boxes, numbers above: \numboxes{from}{to}{per row}{box width cm}{box height cm}
\newcommand{\numboxes}[5]{%
  \begin{tikzpicture}
    \foreach \n in {#1,...,#2} {
      \pgfmathtruncatemacro{\row}{(\n-#1)/#3}
      \pgfmathtruncatemacro{\col}{mod(\n-#1,#3)}
      \pgfmathsetmacro{\x}{\col*#4}
      \pgfmathsetmacro{\y}{-\row*(#5+0.62)}
      \node[anchor=south,inner sep=1pt] at ({\x+#4/2},{\y+0.02}) {\n};
      \draw[line width=0.7pt] (\x,\y) rectangle ({\x+#4},{\y-#5});
    }
  \end{tikzpicture}}
.
\usepackage[letterpaper,top=0.55in,bottom=0.55in,left=0.65in,right=0.65in,
  includeheadfoot,headheight=16pt,headsep=12pt,footskip=26pt]{geometry}
\usepackage{fontspec}
\setmainfont{Andika}
\usepackage{tikz}
\usetikzlibrary{arrows.meta,calc,positioning}
\usepackage{fancyhdr}
\usepackage{array}
\usepackage{enumitem}

\pagestyle{fancy}
\fancyhf{}
\fancyhead[L]{\normalsize Week 7 / Take-away games / \bandname}
\fancyfoot[L]{\normalsize Bellingham Math Circle / Week 7}
\fancyfoot[R]{\normalsize \thepage}
\renewcommand{\headrulewidth}{0.4pt}
\renewcommand{\footrulewidth}{0pt}

\setlength{\parindent}{0pt}
\setlength{\parskip}{0pt}
\raggedright

\definecolor{counterfill}{gray}{0.80}
\definecolor{faded}{gray}{0.62}
\definecolor{lightline}{gray}{0.55}

% \prob{N} starts a problem.
\newcommand{\prob}[1]{\textbf{Problem #1:}\ }

% Ruled writing lines: \writelines{number}{spacing}
\newcommand{\writelines}[2]{%
  \par\begin{tikzpicture}
    \foreach \i in {1,...,#1} {
      \draw[lightline,line width=0.5pt] (0,-\i*#2) -- (\linewidth-1pt,-\i*#2);
    }
    \path (0,0) (0,-#1*#2-0.05);
  \end{tikzpicture}\par}

% One counter at a point: \cnt{(x,y)}{radius}
\newcommand{\cnt}[2]{\draw[line width=0.9pt,fill=counterfill] #1 circle (#2);}
% A counter that the partner already took: dashed, unfilled.
\newcommand{\gone}[2]{\draw[faded,line width=0.9pt,dash pattern=on 2.2pt off 1.8pt] #1 circle (#2);}

% Pile of n counters in rows of 5, top-left counter centre at (x0,y0),
% radius r, pitch p (inside a tikzpicture).  \pileat{n}{x0}{y0}{r}{p}
\newcommand{\pileat}[5]{%
  \foreach \k in {1,...,#1} {
    \pgfmathtruncatemacro{\pilerow}{(\k-1)/5}
    \pgfmathtruncatemacro{\pilecol}{mod(\k-1,5)}
    \pgfmathsetmacro{\pilex}{#2+\pilecol*#5}
    \pgfmathsetmacro{\piley}{#3-\pilerow*#5}
    \cnt{(\pilex,\piley)}{#4}
  }}

% Numbered boxes, numbers above: \numboxes{from}{to}{per row}{box width cm}{box height cm}
\newcommand{\numboxes}[5]{%
  \begin{tikzpicture}
    \foreach \n in {#1,...,#2} {
      \pgfmathtruncatemacro{\row}{(\n-#1)/#3}
      \pgfmathtruncatemacro{\col}{mod(\n-#1,#3)}
      \pgfmathsetmacro{\x}{\col*#4}
      \pgfmathsetmacro{\y}{-\row*(#5+0.62)}
      \node[anchor=south,inner sep=1pt] at ({\x+#4/2},{\y+0.02}) {\n};
      \draw[line width=0.7pt] (\x,\y) rectangle ({\x+#4},{\y-#5});
    }
  \end{tikzpicture}}
